% 第 16 题四联幂函数图（全库首例「选项内嵌四联图」）：源为嵌入位图（2048×481），
% 坐标按源 400 dpi 渲染实测（1 单位 = 1 pt 源页尺寸，x=0.03515cm 即 1pt）。
% 页面坐标（原点左上、y 向下）换算为 y 向上：Y = 400 − y_page。
% 照源要点：四幅坐标系轴长不等、基线不齐平（(D) 幅明显偏右偏低）；曲线逐支按实测点复刻——
% (A) 仅第一象限一支；(B) 关于纵轴对称、过原点（尖点）；(C) 纵轴为渐近线的两支；(D) 奇对称两支、
% 均与原点脱开。选项标注 (A)–(D) 一行平置于全图底部（源实测行心 y=377.6、横向 116.3/216.3/321.6/432.2，
% 非按横轴跨度居中；源内选项行无文字，标注即选项）。
\begin{tikzpicture}[line cap=round,line join=round,font=\small,x=0.03515cm,y=0.03515cm]
% ================= (A) 幅 =================
\draw[->] (113.5,45.9)--(113.5,124.8);   % 纵轴（上端箭头）
\draw[->] (80.8,75.8)--(158.2,75.8);     % 横轴（右端箭头）
\draw plot[smooth] coordinates {(122.8,116.4) (124.6,104.9) (126.7,96.2) (129.6,89.5) (135.4,83.2) (139.7,80.9) (144.0,79.5) (148.3,78.6)};
\node[anchor=center,inner sep=0] at (106.0,68.7) {\(O\)};
\node[anchor=center,inner sep=0] at (119.0,122.0) {\(y\)};
\node[anchor=center,inner sep=0] at (160.5,68.0) {\(x\)};
% ================= (B) 幅 =================
\draw[->] (216.4,45.9)--(216.4,124.8);
\draw[->] (176.0,75.8)--(263.7,75.8);
\draw plot[smooth] coordinates {(179.3,99.6) (188.0,96.6) (196.6,93.2) (205.2,88.8) (213.9,81.9) (216.4,75.8)};
\draw plot[smooth] coordinates {(216.4,75.8) (222.5,85.4) (231.2,90.8) (239.8,94.6) (248.4,97.9) (252.8,99.3)};
\node[anchor=center,inner sep=0] at (209.0,68.7) {\(O\)};
\node[anchor=center,inner sep=0] at (222.0,122.0) {\(y\)};
\node[anchor=center,inner sep=0] at (264.5,68.0) {\(x\)};
% ================= (C) 幅 =================
\draw[->] (321.8,45.2)--(321.8,123.9);
\draw[->] (281.3,74.9)--(369.0,74.9);
\draw plot[smooth] coordinates {(287.4,78.0) (289.9,78.5) (292.2,79.1) (295.3,80.2) (297.7,81.2) (303.8,86.2) (308.1,94.7) (312.4,116.1)};
\draw plot[smooth] coordinates {(330.9,116.4) (331.3,111.5) (331.9,106.8) (332.5,103.4) (333.1,100.2) (333.9,97.8) (335.2,94.2) (336.3,91.5) (337.7,88.8) (339.1,86.6) (340.5,85.0) (342.4,83.3) (346.4,80.7) (351.3,79.0) (354.0,78.3) (356.8,77.6)};
\node[anchor=center,inner sep=0] at (314.4,68.1) {\(O\)};
\node[anchor=center,inner sep=0] at (327.0,121.5) {\(y\)};
\node[anchor=center,inner sep=0] at (370.5,67.0) {\(x\)};
% ================= (D) 幅 =================
\draw[->] (434.6,36.2)--(434.6,123.5);
\draw[->] (387.0,74.6)--(482.0,74.6);
\draw plot[smooth] coordinates {(440.3,116.0) (440.6,114.5) (441.0,112.0) (441.4,110.0) (442.4,104.8) (446.8,93.9) (453.2,87.2) (461.9,83.1) (470.5,81.0) (475.4,80.3)};
\draw plot[smooth] coordinates {(392.7,68.8) (399.2,67.9) (405.7,66.4) (412.2,64.1) (414.6,62.9) (417.0,61.1) (419.4,59.0) (420.9,57.4) (422.1,55.7) (423.0,54.1) (423.7,53.0) (424.5,51.4) (425.2,49.8) (425.7,48.2) (426.2,46.8) (426.7,45.2) (427.2,42.9) (427.7,40.6) (428.2,38.0) (428.5,36.3)};
\node[anchor=center,inner sep=0] at (427.2,67.5) {\(O\)};
\node[anchor=center,inner sep=0] at (440.0,121.1) {\(y\)};
\node[anchor=center,inner sep=0] at (483.5,67.0) {\(x\)};
% ================= 选项标注行（照源：居各幅图正下方） =================
\node[anchor=center,inner sep=0] at (116.3,22.4) {\((A)\)};
\node[anchor=center,inner sep=0] at (216.3,22.4) {\((B)\)};
\node[anchor=center,inner sep=0] at (321.6,22.4) {\((C)\)};
\node[anchor=center,inner sep=0] at (432.2,22.4) {\((D)\)};
\end{tikzpicture}
